\begin{frame}{Récapitulatif des notions vues la semaine dernière}
\centering
    \textbf{\href{https://app.wooclap.com/events/JKAFURZ/questions/6aad0033ad4668ee6c1228cd}{Quiz sur Wooclap}}
\end{frame}

\begin{frame}{Qu'est-ce que le XML?}
\begin{block}{D'abord fut le GML... }
\vspace{1em}
    \textbf{\textit{Generalized Markup Language}}: développé par Charles Goldfarb, Edward Mosher et Raymond Lorie dans les années 1970.\\
    \textrightarrow{} Devient ensuite le Standard Generalized Markup Language (SGML) publié en 1986 et sert de base pour le HTML.\\
    \textrightarrow{} Problème: pas de normes pour utiliser le HTML.
    % Pourquoi ça a été crée?
\end{block}
\vspace{1em}
\begin{block}{... Puis apparut le XML}
\vspace{1em}
\textit{Extensible Markup Language}: publié en 1998 par le W3C (\textit{World Wide Web Consortium})\\ 
\textrightarrow{} Standard international.\\
\textrightarrow{} Sert de base à d'autres langages: XHTML, XLST etc...
\end{block}
    

\end{frame}


\begin{frame}{Pourquoi utiliser le XML?}

    \begin{itemize}
        \item Indiquer la fonction structurelle ou des éléments sémantiques = expliciter certains aspects du texte
        \item Format facile à lire par les machines et l'humain
        \item Interopérable
    \end{itemize}
% Sans ça, c'est nous qui identifions les informations importantes, pas l'ordinateur.
\end{frame}

\begin{frame}{Pourquoi utiliser le XML?}
Qu'est-ce qui peut être relevé dans un texte? \\
\pause 
\begin{itemize}
    \item Type de document: poèmes, essais, pièces de théâtre...
    \item Forme du document: paragraphes, vers, etc...
\end{itemize}
\textbf{Exemples:}
\begin{itemize}
    \item \href{https://lemdo.uvic.ca/qme\_editions/2.0/emdFV\_M.html}{\textit{Queen's Men Editions}}
    \item  \href{https://shelleygodwinarchive.org/sc/oxford/frankenstein/volume/i/\#/p1/mode/xml}{\textit{The Shelley-Godwin Archive}}
\end{itemize}
\end{frame}

